% Part: second-order-logic
% Chapter: metatheory
% Section: introduction

\documentclass[../../../include/open-logic-section]{subfiles}

\begin{document}

\olfileid{sol}{met}{int}

\olsection{परिचय}

प्रथम-क्रम तर्कशास्त्राचे अनेक उपयुक्त गुणधर्म आहेत. ते निर्णेय नाही,
हे आपल्याला माहीत आहे; पण किमान ते स्वयंसिद्धकीकरणीय आहे. म्हणजे
प्रथम-क्रम तर्कशास्त्रासाठी निर्दोष आणि संपूर्ण अशा सिद्धता-प्रणाली
आहेत. त्या असा !!a{derivability} संबंध~$\Proves$ देतात की
!!{sentence}s चा कोणताही संच~$\Gamma$ आणि !!{sentence}~$!A$
यांसाठी $\Gamma \Entails !A$ जर आणि फक्त जर
$\Gamma \Proves !A$. विशेषतः, प्रथम-क्रम तर्कशास्त्रातील
वैधता !!{computably
  enumerable} आहे. एक संगणनक्षम फलन~$f\colon \Nat \to
\Sent[L]$ आहे, ज्यात~$f$ ची मूल्ये नेमकी~$\Lang{L}$ मधील सर्व
वैध !!{sentence}s आहेत. कारण !!{derivation}s चे प्रगणन
करता येते आणि त्यांपैकी एका~!!{sentence} ला !!{derive}
करणाऱ्या निष्पत्तींचे त्या !!{sentence} शी प्रतिचित्रण
करता येते. द्वितीय-क्रम तर्कशास्त्र प्रथम-क्रम तर्कशास्त्रापेक्षा
अधिक अभिव्यक्तिशील आहे; त्यामुळे त्याची वैधता ओळखणे
सर्वसाधारणपणे अधिक कठीण असते. प्रत्यक्षात, द्वितीय-क्रम
तर्कशास्त्र अनिर्णेय तर आहेच, पण त्याची वैधता
!!{computably
  enumerable} सुद्धा नाही, हे आपण दाखवू. म्हणून
द्वितीय-क्रम तर्कशास्त्रासाठी निर्दोष आणि संपूर्ण सिद्धता-प्रणाली
असू शकत नाही (मात्र निर्दोष पण अपूर्ण सिद्धता-प्रणाली उपलब्ध
आहेत आणि त्या संशोधनाचे महत्त्वाचे विषय आहेत).

प्रथम-क्रम तर्कशास्त्राचे आणखी दोन गुणधर्म आहेत: ते संहत
आहे (!!{sentence}s च्या संच~$\Gamma$ चा प्रत्येक सांत उपसंच
पूर्ततायोग्य असेल, तर $\Gamma$ स्वतःही पूर्ततायोग्य असतो)
आणि त्याला लोव्हेनहाइम--स्कोलेम प्रमेय लागू होते
($\Gamma$ चे अनंत प्रतिमान असेल, तर त्याचे
!!a{denumerable} प्रतिमानही असते). हे दोन्ही निष्कर्ष
द्वितीय-क्रम तर्कशास्त्रासाठी लागू होत नाहीत. कारण
द्वितीय-क्रम तर्कशास्त्राला !!{domain}s च्या आकाराविषयी
अशा गोष्टी व्यक्त करता येतात ज्या प्रथम-क्रम तर्कशास्त्राला
व्यक्त करता येत नाहीत.

\end{document}
